%%%%%%%%%%%%%%%%%%%%%%%%
%
% $Autor: Sudhanshu Pukale $
% $Datum: 2025-03-23 $
% $Pfad: D:/Deutschland/Emden-leer/Ml-Attempt2/ML24-EX02-Pico4ML-main/ML24-EX02-Pico4ML-main/Pico4ML-BLE/Contents/General/SPIDisplay.tex $
% $Version: 1 $
%
%%%%%%%%%%%%%%%%%%%%%%%%







\chapter{LCD Display ST7735}

\section{Description}

The LCD display in the Pico4ML-BLE setup for the magic wand project is a compact, color TFT \index{TFT}(Thin-Film Transistor) display module driven by the ST7735 controller, used to provide vibrant visual feedback. Unlike the monochromatic SSD1306 OLED, the ST7735 is a color LCD capable of displaying 16-bit RGB colors (65,536 colors), making it ideal for rendering more complex graphics, text, or sensor data visualizations in real-time. The display has a resolution of 160x80 pixels, which, while relatively low compared to modern screens, is sufficient for the magic wand's purposes, such as displaying spell effects, status indicators, or simple animations. The ST7735 controller manages the display's operation, including pixel rendering and communication over the SPI (Serial Peripheral Interface) protocol.\\



The display's physical size is typically around 1.8 inches diagonally, fitting well within the handheld form factor of the magic wand. It operates at 3.3V, matching the Pico4ML-BLE's logic levels, and consumes more power than an OLED---around 20-40 mA depending on the backlight intensity---due to its TFT technology and backlight requirements. The ST7735 communicates with the Pico4ML-BLE using the SPI protocol, enabling fast data transfer for dynamic updates, which is essential for an interactive device like a magic wand.\cite{Sitronix:2010}



\textbf{Key features include:}

\begin{itemize}

	\item Resolution: 160x80 pixels.

	\item Interface: SPI (4-wire: MOSI, SCK, CS, DC, and often a reset pin; backlight control may require an additional pin).

	\item Power Supply: 3.3V for logic, with a backlight that may require additional current (often controlled via a transistor or PWM).

	\item Size: Typically 1.8 inches diagonally, with a module size of around 34 mm x 54 mm x 3 mm.

	\item Display Type: TFT LCD, offering 16-bit color depth (RGB565 format) and a backlight for visibility in various lighting conditions.\cite{Sitronix:2010}

\end{itemize}



\section{Specifications}

The ST7735 is a widely used controller for small TFT LCD displays, supporting resolutions like 160x80, as seen in 0.96-inch IPS displays. This controller communicates primarily via a 4-wire SPI (Serial Peripheral Interface) protocol, though some variants also support 3-wire SPI or parallel interfaces. For the Pico4ML-BLE, we’ll focus on the 4-wire SPI configuration, which is the most common for this setup.



\begin{itemize}

	\item \textbf{Interface Pins:} The ST7735 requires the following pins for SPI communication:

	\begin{itemize}

		\item MOSI (SDA): Master Out Slave In, for sending data from the Pico4ML-BLE to the display.

		\item SCK (SCL): Serial Clock, for synchronizing data transfer.

		\subitem CS: Chip Select, to enable/disable the display for communication.

		\item DC (or D/C): Data/Command pin, to distinguish between data (pixel information) and commands (e.g., setting the display mode).

		\item RST (Reset): Optional reset pin to initialize the display (can sometimes be tied to VCC if not used).

		\item BLK (Backlight): Controls the backlight; can be connected to a PWM pin for brightness control or left floating for default on (low to turn off).

		\item VCC and GND: Power supply, typically 3.3V for compatibility with the Pico4ML-BLE’s logic levels, though some modules can handle 5V with an internal regulator.\cite{Sitronix:2010}

	\end{itemize}



	

	\item \textbf{Operating Voltage:} The ST7735 typically operates at 3.3V for logic, matching the RP2040’s voltage on the Pico4ML-BLE. The backlight may require 3.3V or 5V depending on the module, often with a current-limiting resistor (e.g., 220 ohms for 5V operation).

	\item \textbf{Color Depth:} The ST7735 supports 16-bit color (65K colors) or 18-bit color (262K colors), though 16-bit is more common for 160x80 displays due to memory constraints on small microcontrollers. Each pixel is represented by 2 bytes in 16-bit mode (5-6-5 RGB format: 5 bits red, 6 bits green, 5 bits blue).\cite{Sitronix:2010}

	\item \textbf{SPI Speed:} The ST7735 can handle SPI clock speeds up to 20 MHz, but for reliable communication with the Pico4ML-BLE, a speed of 1-8 MHz is often used (e.g., 4 MHz is a safe choice).

	\item \textbf{Resolution and Display Type:} The 160x80 resolution means the display has 160 pixels horizontally and 80 pixels vertically, totaling 12,800 pixels. The display is an IPS (In-Plane Switching) TFT LCD, offering wide viewing angles (up to 160 degrees) and vibrant colors compared to older TN or CSTN displays.\cite{Sitronix:2010}

\end{itemize}



\section{Constraints}

\begin{itemize}

	\item SPI Bus Limitations: The Pico4ML-BLE’s RP2040 has two hardware SPI peripherals (SPI0 and SPI1). If other peripherals (e.g., sensors) share the same SPI bus, the CS pin must be carefully managed to avoid conflicts. The ST7735’s SPI communication is unidirectional (master-to-slave), so the MISO pin is not used.

	\item Power Consumption: The ST7735 display consumes around 5-10 mA for the display itself, with the backlight adding another 10-20 mA depending on brightness. For a battery-powered magic wand, this can be a constraint, requiring power-saving techniques like turning off the backlight when not in use or reducing the SPI clock speed to lower power draw.

	\item Memory Requirements: The frame buffer for a 160x80 display in 16-bit color mode requires 160 × 80 × 2 = 25,600 bytes (25.6 KB). The RP2040 has 264 KB of SRAM, so this fits comfortably, but if other tasks (e.g., machine learning models on the Pico4ML-BLE) use significant memory, you may need to optimize by using partial updates or a smaller buffer.

	\item Speed and Refresh Rate: The ST7735 can achieve refresh rates of 30-60 FPS for small displays like this, but the actual rate depends on the SPI speed and the Pico4ML-BLE’s processing. For example, at 4 MHz SPI speed, transferring the full frame buffer (25.6 KB) takes approximately 25,600 × 8 / 4,000,000 = 0.0512 seconds, or about 19.5 FPS. Faster SPI speeds or partial updates can improve this.

	\item Pin Availability: The Pico4ML-BLE has 30 GPIO pins, but using 4-5 pins for the display (MOSI, SCK, CS, DC, optionally RST) reduces the number available for other components like sensors or buttons in the magic wand.\cite{Sitronix:2010}

\end{itemize}



\section{Dimensions}



\begin{itemize}

	\item Physical Size: A typical 0.96-inch 160x80 ST7735 display module measures approximately 24 mm × 30 mm (width × height), with a thickness of about 2-3 mm, including the PCB. The active display area is around 10.8 mm × 21.696 mm, as noted in some 0.96-inch ST7735 modules.

	\item Pin Layout: The module typically has 7-8 pins in a single row (e.g., VCC, GND, CS, RST, DC, SDA/MOSI, SCL/SCK, BLK). The pin pitch is usually 2.54 mm (standard for breadboard compatibility).

	\item Mounting: The display often includes mounting holes or solder pads for integration into a handheld device like a magic wand. Its small size makes it ideal for compact designs.\cite{Sitronix:2010}

\end{itemize}



\section{OS/Software Interfaces/Protocols}

\begin{itemize}

	\item Operating System: The Pico4ML-BLE typically runs bare-metal firmware or a lightweight environment like MicroPython or the Pico SDK (C/C++). There’s no traditional OS, but the firmware handles low-level SPI communication.

	

	\item Software Libraries: Common libraries for the ST7735 include

	\begin{itemize}

		\item MicroPython: adafruit\_st7735 or st7735 libraries, which provide functions for drawing text, shapes, and bitmaps.

		\item C/C++ (Pico SDK): The u8g2 library or Adafruit’s ST7735 library, which offer similar functionality.

		\item Arduino: The Adafruit\_ST7735 library, widely used for Arduino-compatible boards, can be adapted for the Pico4ML-BLE.

	\end{itemize}

	

	\item Protocol: The ST7735 uses a 4-wire SPI protocol. The communication involves:

	\begin{itemize}

		\item Sending commands (e.g., 0x2C to write to the frame buffer) with DC low.

		\item Sending data (e.g., pixel values) with DC high.

		\item The CS pin is pulled low to start communication and high to end it.

	\end{itemize}

	

	\item Initialization: The ST7735 requires a specific initialization sequence, including commands to set the display mode, color format, and orientation. For example, the Adafruit\_ST7735 library typically sends commands like SLPOUT (exit sleep mode), COLMOD (set color mode to 16-bit), and DISPON (turn on display).\cite{Sitronix:2010}

\end{itemize}



\section{Data}

\begin{itemize}

	\item Frame Buffer: The display’s frame buffer is 160 × 80 × 2 = 25.6 KB in 16-bit color mode. Each pixel is a 16-bit value (5-6-5 RGB format), and the entire buffer is sent over SPI to update the display.

	\item Command/Data Structure: The ST7735 distinguishes between commands and data using the DC pin. Commands are 1 byte (e.g., 0x2A to set column address), often followed by data bytes (e.g., start and end column). Pixel data is sent as a continuous stream of 16-bit values.

	\item Orientation: The ST7735 allows the display orientation to be adjusted (horizontal or vertical) via commands like MADCTL (Memory Access Control). For a 160x80 display, the default orientation might be portrait (80x160), but it can be rotated to landscape (160x80) for the magic wand’s user interface.\cite{Sitronix:2010}

\end{itemize}



\section{Quality}



\begin{itemize}

	\item Display Quality: The ST7735-driven 160x80 IPS display offers high contrast and vibrant colors (65K in 16-bit mode). IPS technology provides wide viewing angles (up to 160 degrees), making it readable from various angles, which is useful for a handheld magic wand.

	Refresh and Response: The display has a fast response time (typically 20-30 ms), suitable for simple animations or real-time sensor data display. However, it’s not designed for high-speed video due to SPI bandwidth limitations.

	\item Durability: The display is robust for indoor use, operating in a temperature range of -20°C to 70°C. However, it’s not ruggedized for extreme conditions (e.g., high humidity or shock), so the magic wand design should protect it from physical damage.

	\item Backlight Quality: The backlight provides uniform illumination, but its brightness depends on the current (controlled via the BLK pin). Overdriving the backlight (e.g., with 5V without a resistor) can reduce its lifespan.\cite{Sitronix:2010}

\end{itemize}



\section{Types}

\begin{itemize}

	\item Display Type: This is an IPS TFT LCD, which is superior to older TN or CSTN displays (like the Nokia 6110 LCDs) in terms of color accuracy, contrast, and viewing angles. The ST7735 also supports OLED displays in some configurations, but for 160x80, TFT LCD is more common.

	\item Interface Variants: The ST7735 supports multiple interfaces:

	\begin{itemize}

		\item 4-wire SPI (used here).

		\item 3-wire SPI (combines data and command into one line, using a 9-bit frame).

		\item 8-bit or 16-bit parallel (not practical for the Pico4ML-BLE due to pin constraints).

	\end{itemize}

	\item Color Variants: The display can operate in 12-bit (4-4-4 RGB), 16-bit (5-6-5 RGB), or 18-bit (6-6-6 RGB) color modes, with 16-bit being the most common for this resolution.

\end{itemize}



\section{Structure}

\begin{itemize}

	\item Physical Structure: The display module consists of:

	\begin{itemize}

		\item TFT LCD Panel: The 0.96-inch IPS panel with 160x80 pixels.

		\item ST7735 Controller: Embedded on the module, handling pixel rendering and SPI communication.

		\item PCB: A small breakout board with pin headers for connection.

		\item Backlight: An LED backlight, typically white, with a diffuser for uniform illumination.

		\item FPC (Flexible Printed Circuit): Connects the LCD panel to the PCB.

	\end{itemize}

\end{itemize}



\section{Code with Simple Example}

\begin{code}[h!]
	\captionof{code}{Example LCD Display Code}\label{LCD:example}
	\ArduinoExternal{}{../../Code/Arduino/LCD/LCD.ino}
\end{code}
